La penalización de complejidad de cada árbol se define como:

\begin{equation}
    \Omega(f_k) = \gamma_{\text{xgb}} T_k
    + \frac{1}{2}\lambda_{\text{xgb}}\lVert\mathbf{w}_k\rVert^2.
    \label{eq:xgboost_regularization}
\end{equation}

\noindent\textbf{Donde:}
\begin{itemize}
    \item $\Omega(f_k)$ representa la complejidad penalizada del árbol $k$.
    \item $\gamma_{\text{xgb}}$ representa el costo asociado a añadir hojas al árbol.
    \item $T_k$ representa el número de hojas del árbol $k$.
    \item $\lambda_{\text{xgb}}$ representa el peso de regularización sobre los valores de las hojas.
    \item $\mathbf{w}_k$ representa el vector de pesos o valores asignados a las hojas del árbol $k$.
    \item $\lVert\mathbf{w}_k\rVert^2$ representa la norma cuadrada de dichos pesos.
\end{itemize}

Conceptualmente, esta regularización penaliza tanto la estructura del árbol como la magnitud
de sus predicciones en las hojas. Esto permite que XGBoost mantenga capacidad predictiva sin
crecer de forma innecesaria. En DARL, esta propiedad es útil porque el costo de actualizar o
reentrenar el modelo puede medirse y compararse contra la recuperación de rendimiento.